A basis of a finite Abelian group $\mathfrak{G}$ of order $h$ is a system of elements $B_1, \ldots, B_r$ of orders $h_1, \ldots, h_r$ respectively ($h_i > 1$) such that:

\begin{enumerate}
\item Each element $C \in \mathfrak{G}$ can be represented uniquely as
$$C = B_1^{x_1} B_2^{x_2} \cdots B_r^{x_r},$$
where each $x_i$ runs through a complete residue system modulo $h_i$.
\item Each $h_i$ is a prime power $p^{k_i}$, and $h = h_1 h_2 \cdots h_r$.
\item The orders satisfy $h_1 \geq h_2 \geq \cdots \geq h_r$.
\end{enumerate}

The basis elements are independent. The $h_i$ are called the invariants of $\mathfrak{G}$.
